% GENERATED FILE. Edit canonical lessons, then run npm run generate:books.
% canonical-source-hash: fnv1a64:1eaa6ebb30daa8ef
% canonical-lessons: ML-C185-rajivaykkuka, ML-C185-ariyikkuka, ML-C185-carju-ceyyuka, ML-C185-parakkuka, ML-C185-pintutaruka

\chapter{To resign, To inform, To charge, To fly, To follow}
\label{ch:ml-to-resign-to-inform-to-charge-to-fly-to-follow}
\hlchaptermodality{185}

\begin{chapteropening}
\textbf{By the end of this chapter:} \emph{I can say to resign, to inform, to charge, to fly and to follow.}

Complete the last lesson of chapter 185: I can say to resign, to inform, to charge, to fly and to follow.
\end{chapteropening}

\section[\texorpdfstring{\ml{രാജിവയ്ക്കുക} (rājivaykkuka)}{രാജിവയ്ക്കുക (rājivaykkuka)}]{\ml{രാജിവയ്ക്കുക} (rājivaykkuka) — to resign}

\label{lesson:ML-C185-rajivaykkuka}



\begin{quote}
Before the new one: say the Malayalam for to fill in, then the Malayalam for to sign.
\end{quote}

\subsection*{You'll want to know: \ml{രാജിവയ്ക്കുക}}
\textbf{\ml{രാജിവയ്ക്കുക}} — \emph{rājivaykkuka} — \textquotedblleft{}to resign\textquotedblright{}.

\begin{grammarlens}[title={What you've built}]
One more word for this chapter.
\end{grammarlens}

\subsection*{Your turn}
\begin{itemize}
\raggedright
  \item \emph{Say it:} \emph{rājivaykkuka}
  \item \emph{Say it:} \emph{rājivaykkuka}, once more
  \item \emph{Say it:} say \emph{rājivaykkuka}
  \item \emph{Recall:} say the Malayalam for to fill in, then the Malayalam for to sign, then say \emph{rājivaykkuka} again
\end{itemize}

\begin{tcolorbox}[breakable,colback=teal!4,colframe=teal!35!black,title={Before you move on}]
What does \ml{രാജിവയ്ക്കുക} mean? (\textquotedblleft{}To resign\textquotedblright{}.) Say it once more.
\end{tcolorbox}
\section[\texorpdfstring{\ml{അറിയിക്കുക} (aṟiyikkuka)}{അറിയിക്കുക (aṟiyikkuka)}]{\ml{അറിയിക്കുക} (aṟiyikkuka) — to inform, to let know}

\label{lesson:ML-C185-ariyikkuka}



\begin{quote}
Before the new one: say the Malayalam for to sign, then the Malayalam for to resign.
\end{quote}

\subsection*{You'll want to know: \ml{അറിയിക്കുക}}
\textbf{\ml{അറിയിക്കുക}} — \emph{aṟiyikkuka} — \textquotedblleft{}to inform, to let know\textquotedblright{}.

\begin{grammarlens}[title={What you've built}]
One more word for this chapter.
\end{grammarlens}

\subsection*{Your turn}
\begin{itemize}
\raggedright
  \item \emph{Say it:} \emph{aṟiyikkuka}
  \item \emph{Say it:} \emph{aṟiyikkuka}, once more
  \item \emph{Say it:} say \emph{aṟiyikkuka}
  \item \emph{Recall:} say the Malayalam for to sign, then the Malayalam for to resign, then say \emph{aṟiyikkuka} again
\end{itemize}

\begin{tcolorbox}[breakable,colback=teal!4,colframe=teal!35!black,title={Before you move on}]
What does \ml{അറിയിക്കുക} mean? (\textquotedblleft{}To inform, to let know\textquotedblright{}.) Say it once more.
\end{tcolorbox}
\section[\texorpdfstring{\ml{ചാർജ്} \ml{ചെയ്യുക} (cārjŭ ceyyuka)}{ചാർജ് ചെയ്യുക (cārjŭ ceyyuka)}]{\ml{ചാർജ്} \ml{ചെയ്യുക} (cārjŭ ceyyuka) — to charge (a phone)}

\label{lesson:ML-C185-carju-ceyyuka}



\begin{quote}
Before the new one: say the Malayalam for to resign, then the Malayalam for to inform.
\end{quote}

\subsection*{You'll want to know: \ml{ചാർജ്} \ml{ചെയ്യുക}}
\textbf{\ml{ചാർജ്} \ml{ചെയ്യുക}} — \emph{cārjŭ ceyyuka} — \textquotedblleft{}to charge (a phone)\textquotedblright{}.

\begin{grammarlens}[title={What you've built}]
One more word for this chapter.
\end{grammarlens}

\subsection*{Your turn}
\begin{itemize}
\raggedright
  \item \emph{Say it:} \emph{cārjŭ ceyyuka}
  \item \emph{Say it:} \emph{cārjŭ ceyyuka}, once more
  \item \emph{Say it:} say \emph{cārjŭ ceyyuka}
  \item \emph{Recall:} say the Malayalam for to resign, then the Malayalam for to inform, then say \emph{cārjŭ ceyyuka} again
\end{itemize}

\begin{tcolorbox}[breakable,colback=teal!4,colframe=teal!35!black,title={Before you move on}]
What does \ml{ചാർജ്} \ml{ചെയ്യുക} mean? (\textquotedblleft{}To charge (a phone)\textquotedblright{}.) Say it once more.
\end{tcolorbox}
\section[\texorpdfstring{\ml{പറക്കുക} (paṟakkuka)}{പറക്കുക (paṟakkuka)}]{\ml{പറക്കുക} (paṟakkuka) — to fly}

\label{lesson:ML-C185-parakkuka}



\begin{quote}
Before the new one: say the Malayalam for to inform, then the Malayalam for to charge.
\end{quote}

\subsection*{You'll want to know: \ml{പറക്കുക}}
\textbf{\ml{പറക്കുക}} — \emph{paṟakkuka} — \textquotedblleft{}to fly\textquotedblright{}.

\begin{grammarlens}[title={What you've built}]
One more word for this chapter.
\end{grammarlens}

\subsection*{Your turn}
\begin{itemize}
\raggedright
  \item \emph{Say it:} \emph{paṟakkuka}
  \item \emph{Say it:} \emph{paṟakkuka}, once more
  \item \emph{Say it:} say \emph{paṟakkuka}
  \item \emph{Recall:} say the Malayalam for to inform, then the Malayalam for to charge, then say \emph{paṟakkuka} again
\end{itemize}

\begin{tcolorbox}[breakable,colback=teal!4,colframe=teal!35!black,title={Before you move on}]
What does \ml{പറക്കുക} mean? (\textquotedblleft{}To fly\textquotedblright{}.) Say it once more.
\end{tcolorbox}
\section[\texorpdfstring{\ml{പിന്തുടരുക} (pintuṭaruka)}{പിന്തുടരുക (pintuṭaruka)}]{\ml{പിന്തുടരുക} (pintuṭaruka) — to follow}

\label{lesson:ML-C185-pintutaruka}



\begin{quote}
Before the new one: say the Malayalam for to charge, then the Malayalam for to fly.
\end{quote}

\subsection*{You'll want to know: \ml{പിന്തുടരുക}}
\textbf{\ml{പിന്തുടരുക}} — \emph{pintuṭaruka} — \textquotedblleft{}to follow\textquotedblright{}.

\begin{grammarlens}[title={What you've built}]
One more word for this chapter.
\end{grammarlens}

\subsection*{Your turn}
\begin{itemize}
\raggedright
  \item \emph{Say it:} \emph{pintuṭaruka}
  \item \emph{Say it:} \emph{pintuṭaruka}, once more
  \item \emph{Say it:} say \emph{pintuṭaruka}
  \item \emph{Recall:} say the Malayalam for to charge, then the Malayalam for to fly, then say \emph{pintuṭaruka} again
\end{itemize}

\begin{tcolorbox}[breakable,colback=teal!4,colframe=teal!35!black,title={Before you move on}]
What does \ml{പിന്തുടരുക} mean? (\textquotedblleft{}To follow\textquotedblright{}.) Say it once more.
\end{tcolorbox}
